\practiceweek{3}

\begin{exercisebox}
\begin{pproblem}[$p = 3 \ln n / n$]\label{pp:leader3}
    Redo the analysis with candidate probability $p = \sfrac{3 \ln n}{n}$ and the same $k$.
\end{pproblem}
\tcblower
\begin{wsolution}
    \uses{\scr{thm:leader-atleast-0908}; \scr{lem:disjoint-0908}; \scr{thm:leader-atmost-0908}; \scr{thm:complexity-0908}}
    \leavevmode
    \begin{enumerate}
        \item At least one leader: the candidate with the smallest id always wins at all its
              referees, so there is no leader only if there is no candidate.
              $\prob{\text{no candidate}} = (1-p)^n \leq e^{-pn} = e^{-3 \ln n} = \sfrac{1}{n^3}$.
        \item Two candidates with a common referee cannot both be leaders. Two referee sets are
              disjoint with probability at most $e^{-k^2/n} \leq e^{-\ln n} = \sfrac{1}{n}$, as
              before.
        \item A fixed pair is bad (both are candidates, with no common referee) with probability at
              most $\sfrac{p^2}{n}$. The expected number of bad pairs is at most
              \[\binom{n}{2} \frac{p^2}{n} \leq \frac{n^2}{2} \cdot \frac{9 \ln^2 n}{n^2} \cdot \frac{1}{n} = \frac{9 \ln^2 n}{2n}.\]
        \item By Markov,
              $\prob{\text{more than one leader}} \leq \prob{\text{at least one bad pair}} \leq \frac{9 \ln^2 n}{2n}$.
        \item Messages: at most $2k$ per candidate and $3 \ln n$ candidates in expectation, so at
              most $6k \ln n = \bigO{\sqrt{n} \log^{3/2} n}$.
    \end{enumerate}
    More candidates make ``no leader'' less likely, and ``two leaders'' and the message count
    slightly larger. The asymptotics do not change.
\end{wsolution}
\end{exercisebox}

\begin{exercisebox}
\begin{pproblem}[A triangle with a tail]\label{pp:triangle}
    Vertices $1, 2, 3$ form a triangle and vertex $4$ is joined to vertex $3$.
    \begin{enumerate}[label=(\alph*)]
        \item Write $P$.
        \item Find $\pi$ and verify it.
        \item Does the walk converge?
        \item Starting at vertex $4$, compute $p_1, p_2, p_3$.
    \end{enumerate}
\end{pproblem}
\tcblower
\begin{wsolution}
    \uses{\scr{lem:evolution-0910}; \scr{thm:stationary-0910}; \scr{thm:convergence-0910}}
    (a) The degrees are $2, 2, 3, 1$, so
    \[P = \begin{pmatrix}
            0            & \tfrac{1}{2} & \tfrac{1}{2} & 0            \\
            \tfrac{1}{2} & 0            & \tfrac{1}{2} & 0            \\
            \tfrac{1}{3} & \tfrac{1}{3} & 0            & \tfrac{1}{3} \\
            0            & 0            & 1            & 0
        \end{pmatrix}.\]
    Every row sums to $1$.

    (b) There are $m = 4$ edges, so
    $\pi = \left(\frac{2}{8}, \frac{2}{8}, \frac{3}{8}, \frac{1}{8}\right)
        = \left(\frac{1}{4}, \frac{1}{4}, \frac{3}{8}, \frac{1}{8}\right)$.
    Verify $\pi P = \pi$ entry by entry:
    \begin{itemize}
        \item Vertex $1$: $\pi_2 \cdot \frac{1}{2} + \pi_3 \cdot \frac{1}{3} = \frac{1}{8} + \frac{1}{8} = \frac{1}{4}$.
        \item Vertex $2$: $\pi_1 \cdot \frac{1}{2} + \pi_3 \cdot \frac{1}{3} = \frac{1}{8} + \frac{1}{8} = \frac{1}{4}$.
        \item Vertex $3$: $\pi_1 \cdot \frac{1}{2} + \pi_2 \cdot \frac{1}{2} + \pi_4 \cdot 1
                  = \frac{1}{8} + \frac{1}{8} + \frac{1}{8} = \frac{3}{8}$.
        \item Vertex $4$: $\pi_3 \cdot \frac{1}{3} = \frac{1}{8}$.
    \end{itemize}

    (c) Yes. The graph is connected, and it contains a triangle (an odd cycle), so it is not
    bipartite.

    (d) $p_0 = (0, 0, 0, 1)$.
    \begin{itemize}
        \item $p_1 = (0, 0, 1, 0)$: from $4$ the walk must go to $3$.
        \item $p_2 = \left(\frac{1}{3}, \frac{1}{3}, 0, \frac{1}{3}\right)$: from $3$ each
              neighbour has probability $\frac{1}{3}$.
        \item $p_3 = \left(\frac{1}{3} \cdot \frac{1}{2},\ \frac{1}{3} \cdot \frac{1}{2},\
                  \frac{1}{3} \cdot \frac{1}{2} + \frac{1}{3} \cdot \frac{1}{2} + \frac{1}{3} \cdot 1,\ 0\right)
                  = \left(\frac{1}{6}, \frac{1}{6}, \frac{2}{3}, 0\right)$.
    \end{itemize}
    Each $p_t$ sums to $1$, as it must (a row stochastic matrix preserves probability vectors).
\end{wsolution}
\end{exercisebox}
